\chapter{Complex Numbers and Roots of Unity}

\begin{center}\itshape
Why is $\cos\frac{2\pi}5=\frac{\sqrt5-1}4$, and why does $\sin\frac\pi7\,\sin\frac{2\pi}7\cdots\sin\frac{6\pi}7=\frac7{64}$?\\
Both answers live on the unit circle.
\end{center}

Complex numbers were not invented because someone wanted $\sqrt{-1}$. They were forced on mathematicians, against their will, by an equation
with a perfectly real answer. In 1572 Rafael Bombelli applied Cardano's formula to $x^3=15x+4$, whose root $x=4$ is plain to see, and the formula
produced $\sqrt[3]{2+\sqrt{-121}}+\sqrt[3]{2-\sqrt{-121}}$. Bombelli had the nerve to compute with these ``impossible'' quantities as if they were
numbers, found $(2\pm\sqrt{-1})^3=2\pm\sqrt{-121}$, and got $4$. Two centuries later, an eighteen-year-old Gauss used the same numbers to show that the regular
$17$-gon can be drawn with ruler and compass, something nobody had suspected in two thousand years. The secret, in both cases, is that complex numbers
turn geometry into algebra: a rotation becomes a multiplication, and a regular polygon becomes the set of roots of $z^n=1$.

\section{The exponential form}

We start from the exponential series $\exp z=\sum_{n\ge0}\frac{z^n}{n!}$, which converges absolutely for every $z\in\C$ and satisfies $\exp(a+b)=\exp a\exp b$ (Cauchy product of absolutely convergent series, Chapter~5). We \emph{define} $\cos\theta$ and $\sin\theta$, for $\theta\in\R$, as the real and imaginary parts of $e^{i\theta}=\exp(i\theta)$, and $\pi$ as twice the least positive zero of $\cos$. Splitting the series into even and odd terms gives the usual power series of $\cos$ and $\sin$, hence $\cos'=-\sin$, $\sin'=\cos$, and from the analysis course we use that $\cos$ decreases from $1$ to $-1$ on $[0,\pi]$.

\begin{proposition}[exponential form]\label{prop:expform}
\begin{enumerate}[label=(\alph*)]
\item $|e^{i\theta}|=1$, $e^{i\alpha}e^{i\beta}=e^{i(\alpha+\beta)}$, and hence (de Moivre) $(\cos\theta+i\sin\theta)^n=\cos n\theta+i\sin n\theta$ for $n\in\Z$.
\item Every $z\in\C^*$ can be written $z=r\,e^{i\theta}$ with $r=|z|$ and $\theta\in\R$, unique modulo $2\pi$.
\item (Euler) $\cos\theta=\dfrac{e^{i\theta}+e^{-i\theta}}2$ and $\sin\theta=\dfrac{e^{i\theta}-e^{-i\theta}}{2i}$.
\end{enumerate}
\end{proposition}
\begin{proof}
(a) The series has real coefficients, so $\overline{\exp z}=\exp\overline z$, and $|e^{i\theta}|^2=e^{i\theta}e^{-i\theta}=e^0=1$. The product rule is $\exp(a+b)=\exp a\exp b$, and de Moivre follows by induction (and $e^{-i\theta}=(e^{i\theta})^{-1}$ for negative $n$).

(c) $e^{-i\theta}=\overline{e^{i\theta}}=\cos\theta-i\sin\theta$; add and subtract.

(b) Put $r=|z|$ and $u=z/r=x+iy$, so $x^2+y^2=1$. As $\cos$ maps $[0,\pi]$ onto $[-1,1]$, there is $\theta_0\in[0,\pi]$ with $\cos\theta_0=x$; then $\sin\theta_0^2=1-x^2=y^2$, and $\sin\theta_0\ge0$ on $[0,\pi]$ (it is positive on $]0,\pi[$, since its derivative $\cos$ is positive then negative). So $\sin\theta_0=|y|$, and $\theta=\pm\theta_0$ gives $u=e^{i\theta}$. For uniqueness, $e^{i\theta}=e^{i\theta'}$ means $e^{i(\theta-\theta')}=1$, i.e.\ $\cos(\theta-\theta')=1$; writing $\theta-\theta'=2\pi k+t$ with $t\in[0,2\pi[$ and using $e^{i\pi}=-1$ (since $\cos\pi=-1$ and $|e^{i\pi}|=1$), so that $e^{2i\pi}=1$, we get $\cos t=1$, and on $[0,2\pi[$ this forces $t=0$ (because $\cos t=\cos(2\pi-t)$ and $\cos$ is injective on $[0,\pi]$).
\end{proof}

Two reflexes follow. To \emph{linearise} a power of $\cos$ or $\sin$, expand with Euler and the binomial theorem.
To \emph{factor} a sum like $e^{ia}+e^{ib}$, pull out the half-angle:
\[
\begin{aligned}e^{ia}+e^{ib}&=e^{i\frac{a+b}2}\left(e^{i\frac{a-b}2}+e^{-i\frac{a-b}2}\right)\eqbr=2\cos\frac{a-b}2\;e^{i\frac{a+b}2}.\end{aligned}
\]
In particular $1+e^{i\theta}=2\cos\frac\theta2\,e^{i\theta/2}$ and $1-e^{i\theta}=-2i\sin\frac\theta2\,e^{i\theta/2}$,
so $|1-e^{i\theta}|=2\left|\sin\frac\theta2\right|$.

\begin{example}
\[\begin{aligned}\cos^3\theta&=\frac18\left(e^{i\theta}+e^{-i\theta}\right)^3\eqbr=\frac18\left(e^{3i\theta}+3e^{i\theta}+3e^{-i\theta}+e^{-3i\theta}\right)
\eqbr=\frac{\cos3\theta+3\cos\theta}4.\end{aligned}\]
\end{example}

\section{Roots of unity}

\begin{think}
Draw the solutions of $z^6=1$ in the plane. Which of them have all six as powers, and which only some? What is the sum of all six? Answer from the picture before computing anything.
\end{think}

\begin{proposition}[roots of unity]\label{prop:roots}
Let $n\ge1$ and $\omega=e^{2i\pi/n}$.
\begin{enumerate}[label=(\alph*)]
\item The equation $z^n=1$ has exactly $n$ solutions, the $\omega^k$ for $0\le k\le n-1$. They form a group $\mathbb U_n$ under multiplication, and
$X^n-1=\prod_{k=0}^{n-1}\left(X-\omega^k\right)$.
\item The order of $\omega^k$ is $\frac n{\gcd(k,n)}$. In particular $\omega^k$ is \emph{primitive} (it generates $\mathbb U_n$) exactly when $\gcd(k,n)=1$.
\end{enumerate}
\end{proposition}
\begin{proof}
(a) If $z^n=1$ then $|z|^n=1$, so $|z|=1$ and $z=e^{i\theta}$; then $e^{in\theta}=1$, so $n\theta\in2\pi\Z$ by uniqueness in Proposition~\ref{prop:expform}, i.e.\ $z=\omega^k$ for some $k\in\Z$, and $k$ only matters modulo $n$. Conversely each $\omega^k$ is a solution, and $\omega^0,\dots,\omega^{n-1}$ are distinct (their arguments $\frac{2k\pi}n$ are distinct modulo $2\pi$). Products and inverses of solutions are solutions. Finally, $X^n-1$ has the $n$ distinct roots $\omega^k$, so it is divisible by their product (Chapter~3, Proposition~\ref{prop:rootsfactor}), which is monic of the same degree: they are equal.

(b) Let $d=\gcd(k,n)$. Then $(\omega^k)^m=1$ iff $n\mid km$ iff $\frac nd\mid\frac kd m$ iff $\frac nd\mid m$, the last step by Gauss' lemma since $\gcd\!\left(\frac nd,\frac kd\right)=1$. So the order is $\frac nd$, and $\omega^k$ generates the $n$-element group $\mathbb U_n$ iff its order is $n$, iff $d=1$.
\end{proof}

\begin{figure}[ht]
\centering
\fig{ch02-roots}
\caption{Left: the twelve $12$th roots of unity. The primitive ones, $\omega^k$ with $\gcd(k,12)=1$, are in red; the others generate the smaller groups $\mathbb U_2,\mathbb U_3,\mathbb U_4,\mathbb U_6$. Right: the six chords from $1$ in the regular heptagon have lengths $|1-\omega^k|=2\sin\frac{k\pi}7$, and their product is exactly $7$ (Exercise~\ref{ex:sines}).}
\label{fig:roots}
\end{figure}

\begin{think}
Inscribe a regular heptagon in the unit circle and multiply the lengths of the six chords from one vertex to the others (Figure~\ref{fig:roots}, right). A calculator gives the answer: it is an integer. Which one, and what would you guess for an $n$-gon? The polynomial $X^n-1$ holds the explanation.
\end{think}

\begin{proposition}[orthogonality]\label{prop:orth}
Let $\omega=e^{2i\pi/n}$ and $j\in\Z$. Then
\[
\sum_{k=0}^{n-1}\omega^{jk}=\begin{cases}n&\text{if }n\mid j,\\0&\text{otherwise.}\end{cases}
\]
\end{proposition}
\begin{proof}
If $n\mid j$ every term is $1$. Otherwise $q=\omega^j\ne1$ and $q^n=1$, so the geometric sum is $\frac{q^n-1}{q-1}=0$.
\end{proof}

\section{The roots-of-unity filter}

\begin{think}
How do you get $\sum_k\binom n{2k}$ (the even-index coefficients of $(1+X)^n$) from $(1+1)^n$ and $(1-1)^n$? Now you want only the indices divisible by $3$: what should replace $1$ and $-1$?
\end{think}

Proposition~\ref{prop:orth} lets us extract from a polynomial only the terms whose exponent lies in a given residue class.

\begin{theorem}[roots-of-unity filter]\label{thm:filter}
Let $P(X)=\sum_k a_kX^k$, let $m\ge1$, $\omega=e^{2i\pi/m}$ and $r\in\Z$. Then
\[
\sum_{k\equiv r\ (\mathrm{mod}\ m)}a_k=\frac1m\sum_{j=0}^{m-1}\omega^{-jr}\,P\!\left(\omega^j\right).
\]
\end{theorem}
\begin{proof}
The right-hand side is
\[
\sum_ka_k\cdot\frac1m\sum_{j=0}^{m-1}\omega^{j(k-r)},
\]
and by Proposition~\ref{prop:orth} the inner average is $1$ if $m\mid k-r$ and $0$ otherwise.
\end{proof}

\begin{example}
With $P=(1+X)^n$, $m=2$, $r=0$: $\sum_k\binom n{2k}=\frac12\left(2^n+0^n\right)=2^{n-1}$ for $n\ge1$.
The case $m=3$ is Exercise~\ref{ex:mod3}.
\end{example}

\section{Exercises}

\begin{exercise}[$\star$]\label{ex:lin}
Linearise $\sin^4\theta$, and deduce $\displaystyle\int_0^{\pi}\sin^4\theta\dd\theta$.
\end{exercise}

\begin{exercise}[$\star$]\label{ex:binomcos}
Compute $\displaystyle\sum_{k=0}^n\binom nk\cos(k\theta)$ in closed form.
\end{exercise}

\begin{exercise}[$\star\star$]\label{ex:dirichlet}
\emph{The Dirichlet kernel.} For $\theta\notin2\pi\Z$, show that
\[
D_n(\theta)=\sum_{k=-n}^ne^{ik\theta}=\frac{\sin\left((n+\frac12)\theta\right)}{\sin\frac\theta2},
\]
and deduce $\displaystyle\sum_{k=1}^n\cos k\theta$.
\end{exercise}

\begin{exercise}[$\star\star$]\label{ex:mod3}
How many subsets of $\{1,\dots,n\}$ have a number of elements divisible by $3$?
\end{exercise}

\begin{exercise}[$\star\star$]\label{ex:sines}
Show that $\displaystyle\prod_{k=1}^{n-1}\sin\frac{k\pi}n=\frac n{2^{n-1}}$ for $n\ge2$.
\emph{Hint: evaluate $\prod_{k=1}^{n-1}(X-\omega^k)$ at $X=1$.}
\end{exercise}

\begin{exercise}[$\star\star$]\label{ex:pentagon}
Let $\omega=e^{2i\pi/5}$, $c=\omega+\omega^4$ and $d=\omega^2+\omega^3$. Compute $c+d$ and $cd$, and deduce that
$\cos\frac{2\pi}5=\frac{\sqrt5-1}4$. (This is why the regular pentagon can be drawn with ruler and compass.)
\end{exercise}

\begin{exercise}[$\star\star\star$]\label{ex:circulant}
\emph{Circulant determinants.} Let $a_0,\dots,a_{n-1}\in\C$, $P(X)=\sum_ka_kX^k$, and let $C$ be the $n\times n$ matrix
with entries $C_{r,s}=a_{(s-r)\bmod n}$ (each row is the previous one shifted right by one).
\begin{enumerate}[label=(\alph*)]
\item With $\omega=e^{2i\pi/n}$ and $v_j=\left(1,\omega^j,\omega^{2j},\dots,\omega^{(n-1)j}\right)^{\mathsf T}$, show that $Cv_j=P(\omega^j)\,v_j$.
\item Deduce that $\det C=\prod_{j=0}^{n-1}P(\omega^j)$.
\item Compute the determinant of the circulant matrix with first row $(1,1,0,\dots,0)$.
\end{enumerate}
\end{exercise}

\begin{exercise}[$\star\star\star$]\label{ex:gauss}
\emph{A quadratic Gauss sum.} Let $p$ be an odd prime, $\omega=e^{2i\pi/p}$ and $G=\sum_{a=0}^{p-1}\omega^{a^2}$. Show that $|G|^2=p$.
\end{exercise}

\begin{beauty}
The filter of Theorem~\ref{thm:filter} is the discrete Fourier transform in disguise: averaging $P$ against the
characters $j\mapsto\omega^{-jr}$ isolates one frequency. Tome~4 turns this into the FFT, an algorithm that multiplies
polynomials (and large integers) fast. The Gauss sum of Exercise~\ref{ex:gauss} returns in Tome~1, where its
\emph{square} $G^2=\pm p$ gives a proof of quadratic reciprocity; its sign is a much deeper theorem of Gauss (Tome~4). Exercise~\ref{ex:circulant} is a first example of a matrix
diagonalised by the Fourier basis, the idea behind the characters of Tome~2.
\end{beauty}

\section{Hints}

\begin{hint}{ex:lin}
Expand $\left(\frac{e^{i\theta}-e^{-i\theta}}{2i}\right)^4$ with the binomial formula and group conjugate terms into cosines.
\end{hint}
\begin{hint}{ex:binomcos}
It is the real part of $\sum_k\binom nke^{ik\theta}=(1+e^{i\theta})^n$. Factor out the half-angle: $1+e^{i\theta}=2\cos\frac\theta2\,e^{i\theta/2}$.
\end{hint}
\begin{hint}{ex:dirichlet}
A geometric sum with ratio $e^{i\theta}$. Multiply numerator and denominator by $e^{-i\theta/2}$ to make sines appear.
\end{hint}
\begin{hint}{ex:mod3}
Apply Theorem~\ref{thm:filter} to $(1+X)^n$ with $m=3$, $r=0$. Write $1+j$ and $1+j^2$ in exponential form.
\end{hint}
\begin{hint}{ex:sines}
$\prod_{k=1}^{n-1}(X-\omega^k)=\frac{X^n-1}{X-1}=1+X+\dots+X^{n-1}$. Evaluate at $1$, take moduli, and use $|1-e^{i\theta}|=2\left|\sin\frac\theta2\right|$.
\end{hint}
\begin{hint}{ex:pentagon}
$c+d=\omega+\omega^2+\omega^3+\omega^4=-1$. For $cd$, expand and reduce the exponents modulo $5$. Then $c$ and $d$ are the roots of a quadratic, and $c=2\cos\frac{2\pi}5>0$.
\end{hint}
\begin{hint}{ex:circulant}
(a) The $r$-th coordinate of $Cv_j$ is $\sum_sa_{(s-r)\bmod n}\omega^{sj}$; substitute $t=s-r$ modulo $n$. (b) The $v_j$ form a basis (a Vandermonde matrix). (c) $P=1+X$.
\end{hint}
\begin{hint}{ex:gauss}
$|G|^2=\sum_{a,b}\omega^{a^2-b^2}$. Substitute $a=b+h$: the exponent becomes $h^2+2bh$, linear in $b$. Then sum over $b$ first with Proposition~\ref{prop:orth}.
\end{hint}

\section{Solutions}

\begin{solution}{ex:lin}
\idea ``Linearise'' means: write a power of $\sin$ as a combination of $\cos(k\theta)$, which are easy to integrate. Either expand $\left(\frac{e^{i\theta}-e^{-i\theta}}{2i}\right)^4$ with the binomial formula, or apply the double-angle formula $\sin^2=\frac{1-\cos2\theta}2$ twice, which is quicker here.
\writeup $\sin^2\theta=\frac{1-\cos2\theta}2$, so
\[
\begin{aligned}\sin^4\theta&=\frac{1-2\cos2\theta+\cos^22\theta}4\eqbr=\frac{1-2\cos2\theta+\frac{1+\cos4\theta}2}4\eqbr=\frac{3-4\cos2\theta+\cos4\theta}8.\end{aligned}
\]
The cosines integrate to $0$ over $[0,\pi]$ (for $k\ge1$, $\int_0^\pi\cos 2k\theta\dd\theta=\left[\frac{\sin2k\theta}{2k}\right]_0^\pi=0$), so $\int_0^\pi\sin^4\theta\dd\theta=\frac{3\pi}8$.
\end{solution}

\begin{solution}{ex:binomcos}
\idea A sum $\sum\binom nk\cos(k\theta)$ looks like the binomial formula with $x=e^{i\theta}$, except that $\cos$ is the real part. So compute the complex sum $\sum\binom nke^{ik\theta}=(1+e^{i\theta})^n$, then use the half-angle factorisation to read off its real part.
\writeup The sum is the real part of $\sum_k\binom nke^{ik\theta}=\left(1+e^{i\theta}\right)^n=\left(2\cos\frac\theta2\right)^ne^{in\theta/2}$, so
\[
\sum_{k=0}^n\binom nk\cos(k\theta)=2^n\cos^n\frac\theta2\,\cos\frac{n\theta}2.
\]
\end{solution}

\begin{solution}{ex:dirichlet}
\idea $D_n$ is a geometric sum with ratio $e^{i\theta}\ne1$. The answer we are given contains half-angles, $\sin\frac\theta2$ and $(n+\frac12)\theta$: this is the signal to use the ``pull out the half-angle'' reflex, i.e.\ to make the numerator and denominator symmetric.
\writeup Since $e^{i\theta}\ne1$, the geometric sum gives
\[
D_n(\theta)=e^{-in\theta}\,\frac{e^{i(2n+1)\theta}-1}{e^{i\theta}-1}.
\]
Multiplying numerator and denominator by $e^{-i\theta/2}$,
\[
D_n(\theta)=\frac{e^{i(n+\frac12)\theta}-e^{-i(n+\frac12)\theta}}{e^{i\theta/2}-e^{-i\theta/2}}
=\frac{\sin\left((n+\frac12)\theta\right)}{\sin\frac\theta2}.
\]
Pairing $k$ with $-k$, $D_n(\theta)=1+2\sum_{k=1}^n\cos k\theta$, hence
$\displaystyle\sum_{k=1}^n\cos k\theta=\frac12\left(\frac{\sin\left((n+\frac12)\theta\right)}{\sin\frac\theta2}-1\right)$.
\end{solution}

\begin{solution}{ex:mod3}
\idea Translate: the answer is $\sum_{3\mid k}\binom nk$, the sum of the coefficients of $(1+X)^n$ whose index is in one residue class. That is exactly what the filter of Theorem~\ref{thm:filter} extracts, with $m=3$. The rest is computing $1+j$ in exponential form.
\writeup The number of such subsets is $\sum_{k\equiv0\,(3)}\binom nk$. Apply Theorem~\ref{thm:filter} to $(1+X)^n$ with $m=3$, $r=0$ and
$j=e^{2i\pi/3}$:
\[
\frac13\left(2^n+(1+j)^n+(1+j^2)^n\right).
\]
Now $1+j+j^2=0$ gives $1+j=-j^2=e^{i\pi}e^{-2i\pi/3}=e^{i\pi/3}$ and $1+j^2=\overline{1+j}=e^{-i\pi/3}$, so the answer is $\dfrac{2^n+2\cos\frac{n\pi}3}3$.
For $n=3$ this gives $\frac{8-2}3=2$: the empty set and the whole set.
\end{solution}

\begin{solution}{ex:sines}
\idea A product of sines of angles $\frac{k\pi}n$ suggests the roots of unity, because $|1-e^{i\theta}|=2|\sin\frac\theta2|$ turns $\sin\frac{k\pi}n$ into $\frac12|1-\omega^k|$. So we need $\prod(1-\omega^k)$, which is a value of the polynomial $\prod_{k\ne0}(X-\omega^k)=\frac{X^n-1}{X-1}$.
\writeup Dividing $X^n-1$ by $X-1$ in Proposition~\ref{prop:roots} gives $\prod_{k=1}^{n-1}(X-\omega^k)=1+X+\dots+X^{n-1}$; at $X=1$, $\prod_{k=1}^{n-1}(1-\omega^k)=n$.
Taking moduli, with $|1-e^{i\theta}|=2|\sin\frac\theta2|$ and $\theta=\frac{2k\pi}n$:
\[
n=\prod_{k=1}^{n-1}2\sin\frac{k\pi}n=2^{n-1}\prod_{k=1}^{n-1}\sin\frac{k\pi}n,
\]
all sines being positive since $0<\frac{k\pi}n<\pi$. For $n=7$ this is the $\frac7{64}$ of the opening.
\end{solution}

\begin{solution}{ex:pentagon}
\idea To find two unknown numbers, compute their sum and product: they are then the roots of a quadratic. The relation $1+\omega+\dots+\omega^4=0$ gives the sum, and $\omega^5=1$ lets us reduce exponents in the product. Pairing $\omega$ with $\omega^4=\overline\omega$ makes $c$ real, equal to $2\cos\frac{2\pi}5$.
\writeup By Proposition~\ref{prop:orth}, $1+\omega+\omega^2+\omega^3+\omega^4=0$, so $c+d=-1$. Using $\omega^5=1$,
$cd=\omega^3+\omega^4+\omega^6+\omega^7=\omega^3+\omega^4+\omega+\omega^2=-1$. So $c$ and $d$ are the roots of
$x^2+x-1$, namely $\frac{-1\pm\sqrt5}2$. Since $\omega^4=\overline\omega$, $c=2\cos\frac{2\pi}5>0$ (as $0<\frac{2\pi}5<\frac\pi2$), so
$c=\frac{\sqrt5-1}2$ and $\cos\frac{2\pi}5=\frac{\sqrt5-1}4$.
\end{solution}

\begin{solution}{ex:circulant}
\idea The question hands us candidate eigenvectors; the work in (a) is a change of summation index, which is allowed because indices live modulo $n$ and $\omega^n=1$. For (b), $n$ eigenvectors that form a basis diagonalise $C$, and the determinant is the product of the eigenvalues; the basis property is a Vandermonde determinant. (c) is (b) with $P=1+X$, and $\prod(1+\omega^j)$ is again a value of $X^n-1$, this time at $-1$.
\writeup (a) The $r$-th coordinate of $Cv_j$ is, with $t=s-r$ taken modulo $n$ (allowed since $\omega^{jn}=1$),
\[
\sum_{s=0}^{n-1}a_{(s-r)\bmod n}\,\omega^{js}=\omega^{jr}\sum_{t=0}^{n-1}a_t\,\omega^{jt}=P(\omega^j)\,\omega^{jr}.
\]
(b) The vectors $v_0,\dots,v_{n-1}$ are the columns of a Vandermonde matrix on the distinct numbers $1,\omega,\dots,\omega^{n-1}$; its determinant is non-zero (Chapter~7, Proposition~\ref{prop:vdm}),
so they form a basis. In this basis $C$ is diagonal with entries $P(\omega^j)$, and the determinant, which does not depend on the basis, is their product.

(c) Here $P=1+X$. Evaluating $X^n-1=\prod_j(X-\omega^j)$ at $X=-1$ gives $\prod_j(-1-\omega^j)=(-1)^n-1$, so
\[
\begin{aligned}\det C&=\prod_{j=0}^{n-1}(1+\omega^j)=(-1)^n\left((-1)^n-1\right)\eqbr=1-(-1)^n,\end{aligned}
\]
which is $2$ for $n$ odd and $0$ for $n$ even. (For $n$ even, $(1,-1,1,\dots,-1)$ is in the kernel.)
\end{solution}

\begin{solution}{ex:gauss}
\idea A modulus squared is $G\overline G$, a double sum. Double sums of roots of unity become simple when one variable appears linearly in the exponent, because then orthogonality kills almost everything. The exponent $a^2-b^2$ is quadratic in both, but the substitution $a=b+h$ makes it $h^2+2hb$: linear in $b$.
\writeup Since $\overline\omega=\omega^{-1}$, $|G|^2=G\overline G=\sum_{a,b}\omega^{a^2-b^2}$, where $a,b$ range over $\Z/p\Z$
(the exponent only matters modulo $p$). Substitute $a=b+h$, a bijection for fixed $b$; then $a^2-b^2=h^2+2hb$ and
\[
|G|^2=\sum_{h=0}^{p-1}\omega^{h^2}\sum_{b=0}^{p-1}\omega^{2hb}.
\]
By Proposition~\ref{prop:orth} the inner sum is $p$ if $p\mid2h$ and $0$ otherwise. As $p$ is odd, $p\mid2h$ only for
$h=0$, and $|G|^2=\omega^0\cdot p=p$.
\end{solution}

\section*{Sources and further reading}

\begin{itemize}
\item T.~Needham, \emph{Visual Complex Analysis}, 25th anniversary ed., Oxford University Press, 2023, Chapters~1--2. Complex numbers as geometry, with hundreds of pictures.
\item T.~Andreescu and D.~Andrica, \emph{Complex Numbers from A to\dots Z}, 2nd ed., Birkh\"auser, 2014. Roots of unity and many problems like Exercises~\ref{ex:sines} and~\ref{ex:pentagon}.
\item R.~L.~Graham, D.~E.~Knuth and O.~Patashnik, \emph{Concrete Mathematics}, 2nd ed., Addison-Wesley, 1994, Chapter~7. Generating functions and the roots-of-unity filter.
\item T.~H.~Cormen, C.~E.~Leiserson, R.~L.~Rivest and C.~Stein, \emph{Introduction to Algorithms}, 4th ed., MIT Press, 2022, Chapter~30. Roots of unity at work: the fast Fourier transform and fast polynomial multiplication.
\item B.~C.~Berndt, R.~J.~Evans and K.~S.~Williams, \emph{Gauss and Jacobi Sums}, Wiley, 1998. The full story of the sum $G$ of Exercise~\ref{ex:gauss}, including the determination of its sign.
\end{itemize}
